\section*{Step 266: Shortest Path to RH via Named Externals}

\paragraph{Verdict.}
\[
\mathrm{V\_RH\_attack\_path\_blocked\_by\_external\_gap}.
\]
No active score-3 path was identified.  The best live RH attack interfaces are
score-2 literature-adjacent gates, not published theorems matching the cascade
closure need.

\paragraph{Classification scheme.}
The audit uses four classes:
\[
\begin{array}{ccl}
\text{(i)} & 3 & \text{literature-existing: exact paper statement matches the cascade need},\\
\text{(ii)} & 2 & \text{literature-adjacent: related paper exists, cascade-form theorem absent},\\
\text{(iii)} & 1 & \text{framework-internal claim: need is typed but not in literature},\\
\text{(iv)} & 0 & \text{precision-wall: numerical or conditioning limit}.
\end{array}
\]

\paragraph{Resolved score-3 support.}
Burnol 2002 supplies two exact inherited records.  Theorem 4 gives the Sonine
projection \(\pi_\lambda=P_{L_a^\Gamma}\).  Theorem 8 and equation (1) give
the explicit \(E_\lambda\) and de Branges kernel \(K_a^\Gamma\).  These records
made \(\kappa\) operational at step 201.  They are not active remaining paths;
they expose deeper Branch A/B gates.

\paragraph{Branch A.}
Branch A has three live termini.  The best is the Connes-Consani / adelic
Riemann-Roch recoverability route for the \(M_\zeta\) commutator, classified
as literature-adjacent, score 2.  The relevant Connes-Consani arithmetic-site
and Riemann-Roch papers construct related adelic/arithmetic geometry but do not
state the Sonine Calkin commutator theorem required here.  The faithful Calkin
boundary symbol for \(C^*(P_\infty,M_{m_\ell},P_\eta,I)\) is class iii, score 1.
The closed-form \(\Phi(\sigma,\ell)\) theorem is also class iii, score 1.

\paragraph{Branch B.}
Branch B's best interfaces are Burnol-adjacent.  The transport-sampling theorem
for \(\kappa_i(\tau)\) is class ii, score 2: Burnol 2002/2004 supply adjacent
Sonine/Fourier/zeta machinery, but not the exact sampled boundary identity.
The \(a<1\) linear-combination form needed to resolve CAND1/CAND2 is also class
ii, score 2.  The infinite \(H_\eta\) orthonormal-basis / HS trace theorem is
class iii, score 1.  The inherited \(G^{-1}\) instability is class iv, score 0.

\paragraph{Branch C.}
Branch C is numerically robust but theorem-poor.  The desired theorem saying
that exponential multi-order \(|L_{\rho,k}(G)|\) growth implies a foreclosure
result is class iii, score 1.  Step 252 found no simple zeta- or xi-derivative
bridge, so no literature-existing route is available there.

\paragraph{Ranking.}
The top active interfaces are:
\[
\begin{array}{rllc}
1 & A & \text{Connes-Consani recoverability / adelic bridge} & 2\\
2 & B & \text{Burnol \(a<1\) linear-combination form} & 2\\
3 & B & \text{Burnol transport-sampling theorem} & 2.
\end{array}
\]
Thus the current shortest path is not a direct literature audit; it is the
production of a new cascade-form theorem at one of these score-2 interfaces.
